% Lecture 5 notes: \input by notes.tex, \subimport'ed by the semester book.
% Exam questions (зачёт): 3 (теорема Монтгомери, умножение по модулю без пробных делений),
% 10 (ТЧП, кольца Мерсенна, ТЧП Мерсенна), 11 (ТЧП, кольца Ферма, ТЧП Ферма).
\begin{lecture}{5}{Модулярная арифметика. Метод Монтгомери, кольца Ферма и Мерсенна}

\section{Зачем ускорять модулярную арифметику}

Почти все «законные» операции асимметричной криптографии "--- шифрование и подпись RSA,
согласование ключа Диффи"--~Хеллмана, арифметика на эллиптических кривых "--- сводятся
к длинным цепочкам умножений \emph{по модулю} $m$ длиной в сотни и тысячи бит. Например,
возведение в степень $x^e \bmod m$ с $2048$-битовым показателем требует около трёх тысяч
модулярных умножений (лекция~6). Каждое такое умножение "--- это умножение двух чисел и
\emph{приведение} произведения по модулю $m$, т.\,е. деление с остатком на $m$. Деление
(алгоритм D, лекция~3) "--- самая неудобная операция длинной арифметики: пробные частные,
коррекции, деление двойного слова на слово, которое процессор выполняет в десятки раз дольше
умножения.

В этой лекции разобраны три способа избавиться от деления:
\begin{enumerate}
  \item \emph{метод Монтгомери}: деление на $m$ заменяется делением на степень двойки
    $R = 2^k$, т.\,е. сдвигом; работает для любого нечётного $m$;
  \item \emph{модули специального вида} $2^q - 1$ (числа Мерсенна) и $2^n + 1$ (числа Ферма),
    по которым приведение сводится к сложению или вычитанию двух половин числа;
  \item \emph{теоретико-числовые преобразования} (ТЧП): аналог дискретного преобразования
    Фурье (лекция~4) над кольцом $\Zn{m}$. В кольцах Ферма и Мерсенна корнем из единицы служит
    $2$, и преобразование вычисляется одними сдвигами и сложениями, причём \emph{точно}, без
    ошибок округления.
\end{enumerate}

\section{Арифметика вычетов и её стоимость}

Всюду $m > 1$, вычеты по модулю $m$ хранятся представителями $0 \le x < m$, т.\,е.
$N$-значными числами по основанию $B$, где $N = \floor{\log_B m} + 1$.

\begin{proposition}\label{l5:addmod}
Если $0 \le a, b < m$, то $(a + b) \bmod m$ равно $a + b$ или $a + b - m$, а $(a - b) \bmod m$
равно $a - b$ или $a - b + m$. Сложение и вычитание по модулю стоят $\bigO(N)$.
\end{proposition}

\begin{proof}
$0 \le a + b < 2m$ и $-m < a - b < m$; сравнение и одно сложение или вычитание $N$-значных чисел
стоят $\bigO(N)$ (лекция~2).
\end{proof}

С умножением так не получается: $0 \le ab < m^2$, и частное $\floor{ab/m}$ может быть любым
числом от $0$ до $m - 1$. Стандартный путь "--- умножить ($\Theta(N^2)$ операций школьным методом)
и разделить $2N$-значное число на $N$-значное алгоритмом D (ещё $\Theta(N^2)$ операций, но
с $N$ делениями двойного слова на слово и коррекциями пробных частных). Обратный элемент
$a^{-1} \bmod m$ находится расширенным алгоритмом Евклида (лекция~3) за $\bigO(N^2)$ операций,
поэтому деление по модулю дороже умножения, и в длинных вычислениях его избегают.

\checkyourself{Почему для приведения суммы двух вычетов достаточно одного условного вычитания,
а для произведения "--- нет? Сколько вычитаний $m$ понадобилось бы для приведения $ab$
«в лоб»?}

\section{Метод Монтгомери}

Идея П.~Монтгомери~\cite{montgomery} (1985): заменить деление на «неудобный» модуль $m$
делением на «удобное» число $R$ "--- степень основания системы счисления, т.\,е. сдвиг. Плата
за это "--- лишний множитель $R^{-1}$ в результате, который учитывается специальным
представлением вычетов.

\subsection{Представление Монтгомери}

Пусть $m > 1$ нечётно, $R = 2^k > m$ (обычно $R = B^N$ "--- наименьшая степень основания,
превосходящая $m$). Так как $m$ нечётно, $\gcd(R, m) = 1$, и существуют $R^{-1} \bmod m$ и
\[
  m' = -m^{-1} \bmod R, \qquad\text{т.\,е. } m m' \equiv -1 \pmod R.
\]

\begin{definition}
\emph{Представлением Монтгомери} вычета $x \in \Zn{m}$ называется $\bar x = xR \bmod m$.
\end{definition}

Отображение $x \mapsto \bar x$ "--- биекция $\Zn{m} \to \Zn{m}$ (умножение на обратимый
элемент $R$), причём $\overline{x \pm y} = \bar x \pm \bar y \pmod m$, так что сложение и
вычитание выполняются как обычно. Для произведения
$\bar x \bar y \equiv xy R^2 = \overline{xy}\, R$, т.\,е. после обычного умножения нужно
домножить на $R^{-1}$. Именно эту операцию Монтгомери умеет делать без деления на $m$.

\subsection{Теорема Монтгомери}

\begin{definition}
Для $0 \le T < mR$ положим
\[
  u = \bigl((T \bmod R)\, m'\bigr) \bmod R, \qquad t = \frac{T + um}{R}, \qquad
  \mathrm{REDC}(T) = \begin{cases} t, & t < m,\\ t - m, & t \ge m. \end{cases}
\]
\end{definition}

\begin{theorem}[Монтгомери]\label{l5:redc}
Пусть $m > 1$, $R > m$, $\gcd(R, m) = 1$, $m' = -m^{-1} \bmod R$ и $0 \le T < mR$. Тогда
\begin{enumerate}
  \item $t = (T + um)/R$ "--- целое число;
  \item $\congr{t}{TR^{-1}}{m}$;
  \item $0 \le t < 2m$.
\end{enumerate}
Следовательно, $\mathrm{REDC}(T) = TR^{-1} \bmod m$.
\end{theorem}

\begin{proof}
1. По построению $u \equiv T m' \pmod R$, поэтому
$T + um \equiv T + T m' m \equiv T - T = 0 \pmod R$, т.\,е. $R \mid T + um$.

2. $tR = T + um \equiv T \pmod m$; умножая на $R^{-1} \bmod m$, получаем $t \equiv TR^{-1}$.

3. $t \ge 0$ очевидно, а из $T < mR$ и $u < R$ следует
$T + um < mR + Rm = 2mR$, т.\,е. $t < 2m$.

Итак, $t$ "--- представитель класса $TR^{-1}$ из промежутка $[0, 2m)$; одно условное
вычитание $m$ даёт представителя из $[0, m)$.
\end{proof}

\begin{algorithm}[REDC, редукция Монтгомери]
\begin{algorithmic}[1]
\Require нечётный $m$, $R = 2^k > m$, $m' = -m^{-1} \bmod R$, $0 \le T < mR$.
\Ensure $TR^{-1} \bmod m$.
\State $u \gets \bigl((T \bmod R) \cdot m'\bigr) \bmod R$ \Comment{младшие $k$ бит произведения}
\State $t \gets (T + um) / R$ \Comment{сдвиг на $k$ бит, деление точное}
\If{$t \ge m$} \State $t \gets t - m$ \EndIf
\State \Return $t$
\end{algorithmic}
\end{algorithm}

Все операции здесь "--- умножения, сложения, взятие младших $k$ бит ($\bmod R$) и сдвиг
($/R$); ни одного деления на $m$ «с пробными частными» "--- отсюда название статьи
Монтгомери «Modular multiplication without trial division».

\checkyourself{Где в доказательстве теоремы~\ref{l5:redc} используется условие $T < mR$ и
где "--- $\gcd(R, m) = 1$? Что можно сказать о $t$, если $T < 2mR$?}

\subsection{Умножение по модулю}

\begin{corollary}\label{l5:montmul}
Пусть $0 \le \bar x, \bar y < m$. Тогда $\mathrm{REDC}(\bar x \bar y) = \overline{xy}$.
Кроме того,
\begin{gather*}
  \bar x = \mathrm{REDC}\bigl(x \cdot (R^2 \bmod m)\bigr), \qquad
  x = \mathrm{REDC}(\bar x),\\
  ab \bmod m = \mathrm{REDC}\bigl(\mathrm{REDC}(ab) \cdot (R^2 \bmod m)\bigr).
\end{gather*}
\end{corollary}

\begin{proof}
Во всех случаях аргумент REDC меньше $m^2 < mR$, и теорема~\ref{l5:redc} применима:
$\mathrm{REDC}(\bar x \bar y) \equiv xR\, yR\, R^{-1} = xyR$;
$\mathrm{REDC}(x R^2) \equiv xR$; $\mathrm{REDC}(xR) \equiv x$; наконец,
$\mathrm{REDC}(ab) \equiv abR^{-1}$, и $\mathrm{REDC}(abR^{-1} R^2) \equiv ab$.
\end{proof}

Число $R^2 \bmod m$ вычисляется один раз для данного модуля (одно «настоящее» деление).
Схема вычислений с модулем $m$ такова: перевести входные данные в представление Монтгомери,
выполнить \emph{все} умножения функцией $\mathrm{REDC}(\bar x \bar y)$, сложения и вычитания "---
как обычно, и в конце перевести результат обратно. Для одного умножения выгоды нет (две
редукции вместо одного деления), но для возведения в степень, где умножений тысячи, затраты
на переводы пренебрежимо малы (метод Монтгомери для возведения в степень "--- в лекции~6).

\begin{example}\label{l5:redc-hand}
$m = 97$, $R = 2^7 = 128$. Так как $97 \cdot 33 = 3201 = 25 \cdot 128 + 1$, имеем
$m^{-1} \equiv 33$, $m' = 128 - 33 = 95$; $R^2 = 16384 \equiv 88 \pmod{97}$.
Вычислим $50 \cdot 60 \bmod 97$.

$T = 3000$, $T \bmod 128 = 56$, $u = 56 \cdot 95 \bmod 128 = 5320 \bmod 128 = 72$,
$T + um = 3000 + 72 \cdot 97 = 9984 = 78 \cdot 128$, т.\,е. $\mathrm{REDC}(3000) = 78$.

$T = 78 \cdot 88 = 6864$, $T \bmod 128 = 80$, $u = 80 \cdot 95 \bmod 128 = 48$,
$T + um = 6864 + 4656 = 11520 = 90 \cdot 128$, т.\,е. $\mathrm{REDC}(6864) = 90$.

Итак, $50 \cdot 60 \equiv 90 \pmod{97}$; действительно, $3000 = 30 \cdot 97 + 90$. В
представлении Монтгомери: $\bar a = 50 \cdot 128 \bmod 97 = 95$, $\bar b = 17$,
$\mathrm{REDC}(95 \cdot 17) = \mathrm{REDC}(1615) = 74 = \overline{90}$ и
$\mathrm{REDC}(74) = 90$.
\end{example}

\taskexample{redc-mul}

\subsection{Пословная редукция и стоимость}

В многоразрядной арифметике $R = B^N$, и REDC удобно выполнять по одной цифре: вместо
$m' \bmod B^N$ достаточно одной цифры $m'_0 = -m^{-1} \bmod B$.

\begin{algorithm}[пословная редукция Монтгомери]\label{l5:redc-words}
\begin{algorithmic}[1]
\Require $m = (m_{N-1} \dots m_0)_B$, $\gcd(m, B) = 1$, $R = B^N$,
  $m'_0 = -m^{-1} \bmod B$, $0 \le T < mR$.
\Ensure $TR^{-1} \bmod m$.
\State $t \gets T$
\For{$i = 0, 1, \dots, N - 1$}
  \State $u_i \gets t_i\, m'_0 \bmod B$ \Comment{$t_i$ "--- $i$-я цифра текущего $t$}
  \State $t \gets t + u_i\, m\, B^i$ \Comment{обнуляет $i$-ю цифру}
\EndFor
\State $t \gets t / B^N$ \Comment{сдвиг на $N$ цифр}
\If{$t \ge m$} \State $t \gets t - m$ \EndIf
\State \Return $t$
\end{algorithmic}
\end{algorithm}

\begin{proposition}\label{l5:redc-words-ok}
Алгоритм~\ref{l5:redc-words} вычисляет $TR^{-1} \bmod m$ за $N^2 + N$ умножений цифр.
\end{proposition}

\begin{proof}
Индукцией по $i$: после шага $i$ младшие $i + 1$ цифр $t$ равны нулю. Действительно,
прибавление $u_i m B^i$ не меняет цифр $0, \dots, i - 1$, а $i$-я цифра становится равной
$(t_i + u_i m_0) \bmod B = t_i(1 + m'_0 m_0) \bmod B = 0$, так как $m'_0 m_0 \equiv -1 \pmod B$.
В конце $t = T + Um$, где $U = \sum_i u_i B^i < B^N = R$, и $R \mid t$. Дальше дословно
повторяется доказательство теоремы~\ref{l5:redc} (с $U$ вместо $u$): $t/R \equiv TR^{-1}$ и
$0 \le t/R < 2m$. На шаге $i$ выполняется одно умножение $t_i m'_0$ и $N$ умножений цифр
$u_i m_j$, всего $N(N + 1)$.
\end{proof}

Таким образом, умножение по Монтгомери ($\bar x \bar y$ и редукция) стоит $2N^2 + N$ умножений
цифр "--- столько же, сколько школьное умножение и деление, но без делений двойного слова на
слово, без пробных частных и коррекций. Цикл можно совместить с умножением $\bar x \bar y$
(чередуя прибавление $x_i \bar y B^i$ и редукцию очередной цифры), тогда промежуточный
результат занимает всего $N + 2$ слова. Кроме того, последовательность операций почти не
зависит от данных, что важно для защиты от атак по времени исполнения.

\begin{proposition}\label{l5:nofinal}
Если $4m \le R$ и $0 \le \bar x, \bar y < 2m$, то для $T = \bar x \bar y$ число
$t = (T + um)/R$ из теоремы~\ref{l5:redc} удовлетворяет $0 \le t < 2m$.
Значит, если хранить вычеты в промежутке $[0, 2m)$, финальное вычитание можно не выполнять.
\end{proposition}

\begin{proof}
$T < 4m^2 \le mR$, поэтому $t < (mR + Rm)/R = 2m$; сравнимость $t \equiv TR^{-1}$ доказана в
теореме~\ref{l5:redc} без использования оценки на $T$.
\end{proof}

\subsection{Вычисление \texorpdfstring{$m'$}{m'}}

\begin{proposition}\label{l5:newton}
Пусть $m$ нечётно и $mx \equiv 1 \pmod{2^j}$. Тогда для $x' = x(2 - mx)$ выполнено
$mx' \equiv 1 \pmod{2^{2j}}$.
\end{proposition}

\begin{proof}
$1 - mx' = 1 - 2mx + (mx)^2 = (1 - mx)^2$, а $2^j \mid 1 - mx$, т.\,е. $2^{2j} \mid (1 - mx)^2$.
\end{proof}

Это метод Ньютона для уравнения $1/x - m = 0$ в $2$-адической форме. Так как $m^2 \equiv 1
\pmod 8$ для нечётного $m$, можно начать с $x_0 = m$ (верно по модулю $2^3$): за шаги
$3 \to 6 \to 12 \to 24 \to \dots$ получаем $m^{-1} \bmod 2^k$ за $\bigO(\log k)$ умножений, и
$m' = 2^k - (m^{-1} \bmod 2^k)$. Например, для $m = 97$: $x_0 = 97 \equiv 1 \pmod 8$,
$x_1 = 1 \cdot (2 - 97) = -95 \equiv 33 \pmod{64}$, и уже $97 \cdot 33 \equiv 1 \pmod{128}$.

\subsection{Чётный модуль}

Если $m$ чётно, то $\gcd(R, m) \ne 1$, число $R^{-1} \bmod m$ не существует, и метод
Монтгомери с $R = 2^k$ неприменим (равенство $m m' \equiv -1 \pmod R$ невозможно). Выход "---
китайская теорема об остатках (КТО).

\begin{proposition}\label{l5:even}
Пусть $m = 2^e s$, $s > 1$ нечётно, $x_1 = x \bmod 2^e$, $x_2 = x \bmod s$. Тогда
\[
  x \bmod m = x_1 + 2^e h, \qquad h = \bigl((x_2 - x_1) \cdot (2^{-e} \bmod s)\bigr) \bmod s.
\]
\end{proposition}

\begin{proof}
Число $y = x_1 + 2^e h$ сравнимо с $x_1$ по модулю $2^e$ и
$y \equiv x_1 + 2^e (x_2 - x_1) 2^{-e} = x_2 \pmod s$. Так как $\gcd(2^e, s) = 1$, разность
$x - y$ делится на $2^e s = m$; при этом $0 \le y \le 2^e - 1 + 2^e(s - 1) = m - 1$.
\end{proof}

Для $x = ab$: $x_1 = ab \bmod 2^e$ "--- это просто младшие $e$ бит произведения младших $e$ бит
сомножителей, а $x_2 = ab \bmod s$ вычисляется методом Монтгомери с нечётным модулем $s$
(и $R > s$). Именно так решаются задачи с чётным модулем на семинаре.

\taskexample{redc-mul-even}

\checkyourself{Почему для чётного $m$ метод Монтгомери нельзя «спасти», взяв $R = 3^k$?
Какие операции тогда перестанут быть сдвигами?}

\section{Редукция Барретта}

Другой способ избежать пробных частных предложил П.~Барретт~\cite{barrett} (1986): частное
$\floor{x/m}$ оценивается умножением на заранее вычисленное приближение к $1/m$.
Пусть $m$ "--- $N$-значное число по основанию $B$, $0 \le x < B^{2N}$ и
$\mu = \floor{B^{2N}/m}$ (вычисляется один раз).

\begin{proposition}\label{l5:barrett}
Пусть $q = \floor{x/m}$ и
$\hat q = \Bigl\lfloor \floor{x / B^{N-1}} \cdot \mu \big/ B^{N+1} \Bigr\rfloor$. Тогда
$q - 2 \le \hat q \le q$, т.\,е. $x - \hat q m \in [0, 3m)$, и $x \bmod m$ получается не более
чем двумя вычитаниями $m$.
\end{proposition}

\begin{proof}
Положим $\alpha = x / B^{N-1}$, $\beta = B^{2N}/m$; тогда $q = \floor{\alpha\beta/B^{N+1}}$.
Так как $\floor\alpha \le \alpha$ и $\floor\beta \le \beta$, имеем $\hat q \le q$. С другой
стороны, $\floor\alpha > \alpha - 1$, $\floor\beta > \beta - 1$, а $\alpha < B^{N+1}$ и
$\beta \le B^{N+1}$ (так как $m \ge B^{N-1}$), откуда
\[
  \frac{\floor\alpha \floor\beta}{B^{N+1}} > \frac{\alpha\beta - \alpha - \beta + 1}{B^{N+1}}
  > \frac{\alpha\beta}{B^{N+1}} - 2 \ge q - 2,
\]
и $\hat q \ge q - 2$, поскольку $q - 2$ "--- целое.
\end{proof}

Деления на $B^{N-1}$ и $B^{N+1}$ "--- сдвиги, а в произведениях нужны только старшие (для
$\hat q$) или только младшие (для $x - \hat q m$, вычисляемого по модулю $B^{N+1}$) цифры.
В отличие от метода Монтгомери, редукция Барретта не требует особого представления
вычетов и работает с любым модулем, в том числе чётным.

\begin{example}
$B = 10$, $N = 2$, $m = 97$, $\mu = \floor{10^4/97} = 103$. Для $x = 5000$:
$\hat q = \floor{500 \cdot 103 / 1000} = \floor{51{,}5} = 51$, $x - 51 \cdot 97 = 53$ "---
здесь $\hat q = q$, и вычитания не нужны.
\end{example}

\section{Модули специального вида}

Если модуль можно выбирать (так бывает в криптографических стандартах и при вычислении
свёрток), его выбирают так, чтобы приведение было дешёвым. Классические примеры "--- числа
Мерсенна $2^q - 1$ и числа Ферма $2^n + 1$.

\subsection{Числа Мерсенна}

\begin{proposition}\label{l5:mersenne-red}
Пусть $M = 2^q - 1$ и $x = H \cdot 2^q + L$, $0 \le L < 2^q$. Тогда $\congr{x}{H + L}{M}$.
Если $0 \le x < 2^{2q}$, то не более двух таких «сворачиваний» дают число $y \in [0, M]$,
$y \equiv x$, и $x \bmod M = y$ при $y < M$, $x \bmod M = 0$ при $y = M$.
\end{proposition}

\begin{proof}
$2^q \equiv 1 \pmod M$, поэтому $H \cdot 2^q + L \equiv H + L$. Если $x < 2^{2q}$, то
$H, L \le 2^q - 1$ и $x_1 = H + L \le 2^{q+1} - 2$. Если $x_1 \le M$ "--- готово; иначе
$x_1 = 1 \cdot 2^q + L_1$, $L_1 = x_1 - 2^q \le 2^q - 2$, и $x_2 = 1 + L_1 \le 2^q - 1 = M$.
\end{proof}

Итак, приведение произведения двух $q$-битовых вычетов стоит одного-двух сложений $q$-битовых
чисел "--- $\bigO(N)$ вместо $\Theta(N^2)$ для деления. Аппаратно это сложение с
\emph{циклическим переносом}: перенос из старшего разряда прибавляется к младшему (арифметика
обратного кода). Особенно просто умножение на степень двойки:
$2^j x \bmod M$ "--- \emph{циклический сдвиг} $q$-битовой записи $x$ на $j$ разрядов влево,
так как $2^q \equiv 1$.

\begin{example}
$M = 2^7 - 1 = 127$, $100 \cdot 117 = 11700 = 91 \cdot 2^7 + 52$; $91 + 52 = 143 > 127$,
$143 = 1 \cdot 2^7 + 15$; $1 + 15 = 16$. Итак, $100 \cdot 117 \equiv 16 \pmod{127}$.
\end{example}

\taskexample{mersenne-reduce}

\begin{proposition}\label{l5:mersenne-prime}
Если $2^q - 1$ простое, то $q$ простое.
\end{proposition}

\begin{proof}
Если $q = ab$, $1 < a < q$, то $2^a - 1 \mid (2^a)^b - 1 = 2^q - 1$ и $1 < 2^a - 1 < 2^q - 1$.
\end{proof}

Обратное неверно: $2^{11} - 1 = 2047 = 23 \cdot 89$. Для криптографии важны также
\emph{псевдомерсенновы} модули $m = 2^q - c$ с малым $c$: из $2^q \equiv c$ следует
$H \cdot 2^q + L \equiv L + cH$, и приведение стоит одного умножения на короткое число $c$.
Примеры: простое число Мерсенна $2^{521} - 1$ (кривая NIST P-521) и $2^{255} - 19$
(Curve25519).

\subsection{Числа Ферма}

\begin{proposition}\label{l5:fermat-red}
Пусть $m = 2^n + 1$ и $x = H \cdot 2^n + L$, $0 \le L < 2^n$. Тогда $\congr{x}{L - H}{m}$.
Если $0 \le x \le 2^{2n}$ (например, $x$ "--- произведение двух вычетов), то
$x \bmod m = L - H$ при $L \ge H$ и $L - H + m$ при $L < H$.
\end{proposition}

\begin{proof}
$2^n \equiv -1 \pmod m$. При $x \le 2^{2n}$ имеем $H \le 2^n$, поэтому
$-2^n \le L - H \le 2^n - 1$, т.\,е. $L - H \in (-m, m)$.
\end{proof}

\begin{example}
$m = 2^8 + 1 = 257$, $200 \cdot 150 = 30000 = 117 \cdot 2^8 + 48$, $48 - 117 = -69$,
$-69 + 257 = 188$. Итак, $200 \cdot 150 \equiv 188 \pmod{257}$.
\end{example}

\taskexample{fermat-reduce}

\begin{proposition}\label{l5:fermat-prime}
Если $2^n + 1$ простое, то $n$ "--- степень двойки.
\end{proposition}

\begin{proof}
Если $n = kl$ с нечётным $k > 1$, то $2^l + 1 \mid (2^l)^k + 1 = 2^n + 1$ (так как
$x + 1 \mid x^k + 1$ при нечётном $k$), и $1 < 2^l + 1 < 2^n + 1$.
\end{proof}

Числа $F_t = 2^{2^t} + 1$ называются \emph{числами Ферма}. Числа $F_0, \dots, F_4 = 3, 5, 17,
257, 65537$ простые, но уже $F_5 = 2^{32} + 1 = 641 \cdot 6700417$ (Эйлер); других простых
чисел Ферма не известно. Для ТЧП (см. ниже) простота модуля не нужна.

\subsection{Код с уменьшением на единицу}

Вычеты по модулю $2^n + 1$ "--- это $0, 1, \dots, 2^n$, т.\,е. $2^n + 1$ значение, которые
не помещаются в $n$ бит; лишний $(n{+}1)$-й бит нужен ровно одному значению
$2^n \equiv -1$. Л.~Лейбовиц~\cite{leibowitz} предложил хранить вместо $x$ число $x - 1$.

\begin{definition}
\emph{Кодом с уменьшением на единицу} (diminished-one) вычета $x \in \{0, 1, \dots, 2^n\}$
называется $(n{+}1)$-битовое слово: при $x \ne 0$ "--- старший бит $0$ и младшие $n$ бит
$x' = x - 1$; при $x = 0$ "--- старший бит (признак нуля) $1$ и нули.
\end{definition}

Так, при $n = 8$: $1 \mapsto 0\,00000000$, $256 \equiv -1 \mapsto 0\,11111111$,
$0 \mapsto 1\,00000000$. Нуль обрабатывается отдельно по признаку (сумма с нулём и
произведение на нуль тривиальны), а для ненулевых вычетов все операции выполняются над
$n$-битовыми словами $x'$.

\begin{proposition}\label{l5:dim1}
Пусть $m = 2^n + 1$, $x, y \not\equiv 0$, $x' = x - 1$, $y' = y - 1 \in [0, 2^n)$.
\begin{enumerate}
  \item (отрицание) $(-x)' = 2^n - 1 - x'$ "--- побитовое дополнение $x'$;
  \item (сложение) пусть $s = x' + y'$; если $s \ge 2^n$ (есть перенос), то
    $(x + y)' = s - 2^n$; если $s = 2^n - 1$, то $x + y \equiv 0$; иначе $(x + y)' = s + 1$;
  \item (умножение на $2$) $(2x)' = (2x' \bmod 2^n) + (1 - c)$, где $c$ "--- старший бит $x'$:
    сдвиг влево, в младший разряд записывается инвертированный выдвинутый бит;
  \item (умножение) пусть $P = x'y' + x' + y' = H \cdot 2^n + L$; тогда
    $(xy)' \equiv L - H \pmod m$, и $xy \equiv 0$, только если $(L - H) \bmod m = 2^n$
    (для простого $m$ это невозможно).
\end{enumerate}
\end{proposition}

\begin{proof}
1. $-x \equiv 2^n + 1 - x \in [1, 2^n]$, поэтому $(-x)' = 2^n - x = 2^n - 1 - x'$.

2. $x + y - 1 = s + 1$. При $s \ge 2^n$ получаем $s + 1 - m = s - 2^n \in [0, 2^n - 2]$. При
$s < 2^n$ число $s + 1 \in [1, 2^n]$; оно равно $2^n$, только если $x + y = 2^n + 1 \equiv 0$.

3. $2x - 1 = 2x' + 1$. При $c = 0$ это число меньше $2^n$; при $c = 1$ вычитаем $m$:
$2x' + 1 - 2^n - 1 = 2x' - 2^n$. Нуля не возникает, так как $m$ нечётно.

4. $xy - 1 = (x' + 1)(y' + 1) - 1 = P \equiv L - H$ по предложению~\ref{l5:fermat-red}.
Код ненулевого вычета "--- число из $[0, 2^n)$; значение $2^n$ отвечает $xy - 1 \equiv -1$,
т.\,е. $xy \equiv 0$, что для простого $m$ при $x, y \not\equiv 0$ невозможно.
\end{proof}

\taskexample{dim1-code}

\checkyourself{Запишите коды с уменьшением на единицу чисел $0$, $1$, $2$, $-1$, $-2$ по
модулю $257$. Чему равен код суммы $128 + 129$?}

\section{Теоретико-числовые преобразования}

В лекции~4 свёртка последовательностей вычислялась с помощью дискретного преобразования Фурье
(ДПФ) над $\CC$ с корнем из единицы $e^{2\pi i/n}$. Вычисления в плавающей точке дают ошибки
округления, и для точного результата приходится следить за точностью. Заметим, что
доказательства формулы обращения ДПФ и теоремы о свёртке используют лишь два свойства корня
$\omega$: $\omega^n = 1$ и $\sum_{j} \omega^{jk} = 0$ при $k \not\equiv 0 \pmod n$. Если найти
такой элемент в кольце вычетов $\Zn{m}$, все вычисления станут целочисленными и точными.

\subsection{Главные корни из единицы}

Пусть $A$ "--- коммутативное кольцо с единицей (у нас $A = \Zn{m}$), $n \ge 2$.

\begin{definition}
Элемент $\omega \in A$ называется \emph{главным корнем степени $n$ из единицы}, если
$\omega^n = 1$ и
\[
  \sum_{j=0}^{n-1} \omega^{jk} = 0 \qquad \text{при } k = 1, 2, \dots, n - 1.
\]
\end{definition}

\begin{lemma}\label{l5:principal}
Если $\omega^n = 1$ и элементы $\omega^k - 1$ обратимы при $1 \le k \le n - 1$, то $\omega$ "---
главный корень степени $n$ из единицы. В поле любой элемент порядка ровно $n$ "--- главный
корень степени $n$.
\end{lemma}

\begin{proof}
$(\omega^k - 1) \sum_{j=0}^{n-1} \omega^{jk} = \omega^{nk} - 1 = 0$; умножая на
$(\omega^k - 1)^{-1}$, получаем, что сумма равна нулю. В поле $\omega^k - 1 \ne 0$ обратимо при
$1 \le k < n$, если порядок $\omega$ равен $n$.
\end{proof}

\begin{theorem}\label{l5:roots}
Пусть $m > 1$, $\omega \in \Zn{m}$, $\omega^n \equiv 1 \pmod m$. Следующие условия
равносильны:
\begin{enumerate}
  \item $\omega$ "--- главный корень степени $n$ из единицы и $\gcd(n, m) = 1$;
  \item для каждого простого $p \mid m$ порядок $\omega$ по модулю $p$ равен $n$.
\end{enumerate}
В этом случае $n \mid p - 1$ для всех простых $p \mid m$.
\end{theorem}

\begin{proof}
$(2) \Rightarrow (1)$. При $1 \le k < n$ имеем $\omega^k \not\equiv 1 \pmod p$ для каждого
$p \mid m$, т.\,е. $\gcd(\omega^k - 1, m) = 1$ и $\omega^k - 1$ обратим по модулю $m$. По
лемме~\ref{l5:principal} $\omega$ "--- главный корень. Порядок элемента группы $\Znx{p}$ делит
её порядок $p - 1$, поэтому $n \mid p - 1$, $p \nmid n$ для всех $p \mid m$, и $\gcd(n, m) = 1$.

$(1) \Rightarrow (2)$. Пусть $p \mid m$, $d$ "--- порядок $\omega$ по модулю $p$; так как
$\omega^n \equiv 1$, имеем $d \mid n$. Если $d < n$, то при $k = d$ все слагаемые
$\omega^{jd} \equiv 1 \pmod p$ и $\sum_{j=0}^{n-1} \omega^{jd} \equiv n \pmod p$. Но эта сумма
сравнима с нулём по модулю $m$, а значит, и по модулю $p$, т.\,е. $p \mid n$ "--- противоречие с
$\gcd(n, m) = 1$. Итак, $d = n$.
\end{proof}

\begin{corollary}\label{l5:roots-prime}
По простому модулю $p$ главный корень степени $n$ из единицы существует тогда и только тогда,
когда $n \mid p - 1$.
\end{corollary}

\begin{proof}
Необходимость "--- теорема~\ref{l5:roots}. Достаточность: группа $\Znx{p}$ циклическая
(существует первообразный корень $g$, лекция~7), и $\omega = g^{(p-1)/n}$ имеет порядок $n$.
\end{proof}

Например, $p = 998244353 = 119 \cdot 2^{23} + 1$ "--- простое, $3$ "--- первообразный корень,
и $3^{119}$ "--- главный корень степени $2^{23}$: по такому модулю можно вычислять свёртки
длины до $2^{23}$ алгоритмом БПФ. Но умножения в нём "--- общие, с приведением по модулю $p$.
Кольца Мерсенна и Ферма интересны тем, что в них корнем служит $\pm 2$ или степень двойки.

\subsection{Преобразование и теорема о свёртке}

\begin{definition}
Пусть $\omega \in \Zn{m}$ "--- главный корень степени $n$ из единицы, $\gcd(n, m) = 1$.
\emph{Теоретико-числовым преобразованием} (ТЧП) последовательности
$a = (a_0, \dots, a_{n-1})$, $a_j \in \Zn{m}$, называется последовательность
\[
  A_k = \sum_{j=0}^{n-1} a_j \omega^{jk} \bmod m, \qquad k = 0, 1, \dots, n - 1.
\]
\emph{Циклической свёрткой} последовательностей $a$ и $b$ длины $n$ называется
$c = a * b$, $c_j = \sum_{i + l \equiv j \ (\mathrm{mod}\ n)} a_i b_l$.
\end{definition}

\begin{theorem}\label{l5:inverse}
В условиях определения
\[
  a_j = n^{-1} \sum_{k=0}^{n-1} A_k \omega^{-jk} \bmod m, \qquad j = 0, \dots, n - 1,
\]
т.\,е. обратное преобразование "--- это ТЧП с корнем $\omega^{-1} = \omega^{n-1}$, умноженное
на $n^{-1} \bmod m$.
\end{theorem}

\begin{proof}
Сначала \emph{соотношение ортогональности}: $\sum_{k=0}^{n-1} \omega^{rk} = n$ при
$r \equiv 0 \pmod n$ и $0$ иначе. Первое очевидно ($\omega^n = 1$), второе "--- определение
главного корня для $r \bmod n \in \{1, \dots, n-1\}$ (так как $\omega^{rk}$ зависит лишь от
$r \bmod n$). Теперь
\[
  \sum_{k} A_k \omega^{-jk} = \sum_k \sum_l a_l \omega^{(l - j)k}
  = \sum_l a_l \sum_k \omega^{(l-j)k} = n a_j,
\]
так как внутренняя сумма отлична от нуля только при $l = j$.
\end{proof}

\begin{theorem}[о свёртке]\label{l5:convolution}
ТЧП переводит циклическую свёртку в поточечное произведение: если $c = a * b$, то
$C_k = A_k B_k$ для всех $k$.
\end{theorem}

\begin{proof}
Так как $\omega^n = 1$, значение $\omega^{jk}$ зависит лишь от $j \bmod n$, поэтому
\[
  C_k = \sum_j \sum_{i + l \equiv j} a_i b_l\, \omega^{jk}
  = \sum_{i, l} a_i b_l\, \omega^{(i + l)k}
  = \Bigl(\sum_i a_i \omega^{ik}\Bigr)\Bigl(\sum_l b_l \omega^{lk}\Bigr) = A_k B_k. \qedhere
\]
\end{proof}

\begin{algorithm}[свёртка с помощью ТЧП]
\begin{algorithmic}[1]
\Require $a, b \in (\Zn{m})^n$; главный корень $\omega$ степени $n$, $\gcd(n, m) = 1$.
\Ensure $c = a * b \bmod m$.
\State $A \gets \mathrm{NTT}_\omega(a)$, $B \gets \mathrm{NTT}_\omega(b)$
\State $C_k \gets A_k B_k \bmod m$, $k = 0, \dots, n - 1$
\State $c \gets n^{-1} \cdot \mathrm{NTT}_{\omega^{-1}}(C) \bmod m$
\end{algorithmic}
\end{algorithm}

По определению каждое ТЧП стоит $n^2$ умножений в $\Zn{m}$; при $n = 2^s$ алгоритм БПФ
(лекция~4) переносится на ТЧП дословно (он использует только $\omega^{n/2} = -1$ и
$\omega^n = 1$) и даёт $\bigO(n \log n)$ операций в $\Zn{m}$.

\begin{proposition}\label{l5:exact}
Пусть $0 \le a_i, b_i \le W - 1$ "--- целые числа и $m > n (W - 1)^2$. Тогда результат
свёртки с помощью ТЧП (вычеты из $[0, m)$) совпадает с циклической свёрткой над $\ZZ$. Если
$a$ и $b$ дополнены нулями так, что $a_i = 0$ при $i \ge L_1$, $b_i = 0$ при $i \ge L_2$ и
$n \ge L_1 + L_2 - 1$, то циклическая свёртка совпадает с линейной.
\end{proposition}

\begin{proof}
Каждое $c_j$ над $\ZZ$ "--- сумма не более $n$ произведений, не превосходящих $(W - 1)^2$, т.\,е.
$0 \le c_j < m$, и вычет $c_j \bmod m$ равен самому $c_j$. Во втором случае при $i < L_1$,
$l < L_2$ имеем $i + l \le L_1 + L_2 - 2 < n$, так что «заворачивания» индексов по модулю $n$
не происходит.
\end{proof}

Таким образом, умножение чисел через ТЧП устроено так же, как через БПФ (лекция~4): цифры
по основанию $W$ "--- это последовательности, их линейная свёртка "--- коэффициенты
произведения, после чего выполняются переносы.

\begin{example}\label{l5:ntt17}
Умножим $23 \cdot 14$ по модулю $m = 17 = 2^4 + 1$ с $n = 4$, $\omega = 4$
($4^2 = 16 \equiv -1$, $4^4 \equiv 1$). Цифры: $a = (3, 2, 0, 0)$, $b = (4, 1, 0, 0)$;
оценка предложения~\ref{l5:exact} здесь не помогает ($2 \cdot 9^2 > 17$), но коэффициенты
линейной свёртки $(12, 11, 2, 0)$ меньше $17$, так что результат будет точным.
\begin{gather*}
  A_k = 3 + 2 \cdot 4^k:\ A = (5, 11, 1, 12), \qquad
  B_k = 4 + 4^k:\ B = (5, 8, 3, 0),\\
  C = (25, 88, 3, 0) \bmod 17 = (8, 3, 3, 0).
\end{gather*}
Обратное: $\omega^{-1} = 13 \equiv -4$, $n^{-1} = 4^{-1} = 13 \equiv -4$;
$\sum_k C_k \omega^{-jk} = (14, -7, 8, 17) \equiv (14, 10, 8, 0)$, и умножение на $13$ даёт
$c = (12, 11, 2, 0)$. Итак, $23 \cdot 14 = 12 + 11 \cdot 10 + 2 \cdot 100 = 322$.
\end{example}

\checkyourself{Почему в примере~\ref{l5:ntt17} нельзя взять $n = 2$? Какое максимальное
значение цифры $W - 1$ гарантирует точный результат при $n = 4$ и $m = 257$?}

\subsection{Теоретико-числовое преобразование Мерсенна}

\begin{theorem}[Рейдер~\cite{rader-mersenne}]\label{l5:mnt}
Пусть $q$ "--- нечётное простое, $M = 2^q - 1$ (не обязательно простое). Тогда
\begin{enumerate}
  \item $\omega = 2$ "--- главный корень степени $q$ из единицы в $\Zn{M}$;
  \item $\omega = -2$ "--- главный корень степени $2q$ из единицы в $\Zn{M}$;
\end{enumerate}
и числа $q$, $2q$ обратимы по модулю $M$.
\end{theorem}

\begin{proof}
Применим теорему~\ref{l5:roots}. Пусть $p \mid M$ "--- простое; $p$ нечётно.

1. $2^q = M + 1 \equiv 1 \pmod p$, поэтому порядок $2$ по модулю $p$ делит простое $q$; он не
равен $1$, так как $2 \not\equiv 1 \pmod p$. Значит, он равен $q$.

2. $(-2)^{2q} = 4^q \equiv 1$, поэтому порядок $d$ числа $-2$ делит $2q$, т.\,е.
$d \in \{1, 2, q, 2q\}$. Случаи $d = 1, 2$ означают $p \mid 3$ ($-2 \equiv 1$ или $4 \equiv 1$),
но $3 \nmid M$: для нечётного $q$ имеем $2^q \equiv (-1)^q = -1 \pmod 3$ и
$M \equiv -2 \equiv 1 \pmod 3$. Случай $d = q$ невозможен: $(-2)^q = -2^q \equiv -1
\not\equiv 1 \pmod p$. Итак, $d = 2q$.

Обратимость $q$ и $2q$ следует из теоремы~\ref{l5:roots}.
\end{proof}

\emph{Преобразованием Мерсенна} называется ТЧП по модулю $M = 2^q - 1$ с $\omega = \pm 2$.
Умножение на $\omega^{jk} = 2^{jk \bmod q}$ "--- циклический сдвиг $q$-битового слова, поэтому
прямое и обратное преобразования не требуют умножений вовсе: только сдвиги и сложения с
циклическим переносом. Умножения нужны лишь для поточечного произведения $C_k = A_k B_k$
(с дешёвым приведением по модулю $M$) и для домножения на $n^{-1}$.

Недостаток: длина преобразования $n = q$ или $2q$ "--- простое число или удвоенное простое,
поэтому быстрый алгоритм БПФ по основанию $2$ неприменим, и само преобразование стоит
$\Theta(n^2)$ сложений. Кроме того, длина жёстко связана с длиной слова: для $32$-битовой
арифметики $M = 2^{31} - 1$ и $n \le 62$.

\begin{example}
$M = 2^{11} - 1 = 2047 = 23 \cdot 89$ "--- составное, но порядок $2$ по модулю и $23$, и $89$
равен $11$, а порядок $-2$ "--- $22$. Поэтому в $\Zn{2047}$ есть ТЧП длины $11$ и $22$.
\end{example}

\taskexample{ntt-mersenne}

\subsection{Теоретико-числовое преобразование Ферма}

\begin{theorem}[Агарвал"--~Буррас~\cite{agarwal-burrus}]\label{l5:fnt}
Пусть $F_t = 2^b + 1$, $b = 2^t$, $t \ge 1$. Тогда в $\Zn{F_t}$
\begin{enumerate}
  \item $\omega = 2$ "--- главный корень степени $2b = 2^{t+1}$ из единицы, и для любого
    $n = 2^s \le 2b$ элемент $2^{2b/n}$ "--- главный корень степени $n$;
  \item при $t \ge 2$ элемент $\omega = 2^{b/4}(2^{b/2} - 1)$ удовлетворяет
    $\omega^2 \equiv 2$ и является главным корнем степени $4b$;
  \item $n^{-1} \equiv -2^{b - s}$ для $n = 2^s \le 2b$.
\end{enumerate}
\end{theorem}

\begin{proof}
Пусть $p \mid F_t$ "--- простое; $p$ нечётно, и $2^b \equiv -1 \pmod p$.

1. $2^{2b} \equiv 1$, поэтому порядок $2$ по модулю $p$ "--- делитель $2^{t+1}$, т.\,е. степень
двойки. Если бы он делил $2^t = b$, то $2^b \equiv 1$ и $-1 \equiv 1 \pmod p$, т.\,е. $p = 2$.
Значит, порядок равен $2b$, и по теореме~\ref{l5:roots} $2$ "--- главный корень степени $2b$.
Порядок $2^{2b/n}$ по модулю каждого $p$ равен $2b / (2b/n) = n$.

2. Используя $2^b \equiv -1$:
$\omega^2 = 2^{b/2}(2^b - 2^{b/2+1} + 1) = 2^{3b/2} - 2^{b+1} + 2^{b/2}
\equiv -2^{b/2} + 2 + 2^{b/2} = 2$. Тогда $\omega^{4b} = 2^{2b} \equiv 1$ и
$\omega^{2b} = 2^b \equiv -1 \not\equiv 1 \pmod p$, поэтому порядок $\omega$ по модулю $p$ "---
степень двойки, делящая $4b$, но не $2b$, т.\,е. $4b$.

3. $2^s \cdot (-2^{b-s}) = -2^b \equiv 1$.
\end{proof}

\emph{Преобразование Ферма} (FNT) "--- это ТЧП по модулю $F_t$ с $\omega = 2^{2b/n}$. Его
достоинства:
\begin{itemize}
  \item длина $n$ "--- степень двойки, поэтому применим алгоритм БПФ: $\bigO(n \log n)$
    операций, причём все «бабочки» $(x, y) \mapsto (x + \omega^k y,\ x - \omega^k y)$ сводятся
    к сдвигу и сложению/вычитанию с приведением по предложению~\ref{l5:fermat-red};
  \item $n^{-1} = -2^{b-s}$ "--- тоже сдвиг;
  \item для хранения вычетов достаточно $b$-битовых слов в коде с уменьшением на единицу.
\end{itemize}
Недостаток "--- та же жёсткая связь длины с размером слова: $n \le 2b$ (или $4b$ с
корнем $\sqrt 2$). Например, по модулю $F_2 = 17$ корень $2$ имеет порядок $8$,
по модулю $F_3 = 257$ "--- порядок $16$, а элемент $2^2(2^4 - 1) = 60$ "--- порядок $32$
($60^2 = 3600 \equiv 2 \pmod{257}$).

\taskexample{ntt-fermat}

\checkyourself{Какой корень $\omega$ (степень двойки) нужен для ТЧП длины $8$ по модулю $257$?
Чему равны $\omega^{-1}$ и $n^{-1}$? Какие из этих умножений "--- сдвиги?}

\subsection{Ограничения и выбор параметров}

Пусть требуется точная линейная свёртка двух последовательностей длины $L$ из $w$-битовых
цифр ($W = 2^w$). Нужно $n \ge 2L - 1$ и $m > L(W - 1)^2$ (ср. предложение~\ref{l5:exact}),
т.\,е. $\log_2 m \gtrsim 2w + \log_2 L$. Для колец Ферма и Мерсенна длина преобразования
ограничена размером модуля (каждый коэффициент линейной свёртки "--- сумма не более $L$
произведений цифр):
\begin{itemize}
  \item $F_5 = 2^{32} + 1$: $n \le 2b = 64$ с $\omega = 2$ (или $128$ с $\omega = \sqrt 2$);
    при $n = 128$, $L = 64$ годятся $13$-битовые цифры ($64 \cdot (2^{13} - 1)^2 < 2^{32}$),
    т.\,е. перемножаются числа длиной $64 \cdot 13 = 832$ бита;
  \item $M_{31} = 2^{31} - 1$: $n \le 62$ с $\omega = -2$, $L = 31$, тоже $13$-битовые цифры
    ($31 \cdot (2^{13} - 1)^2 < 2^{31} - 1$), т.\,е. числа длиной $403$ бита.
\end{itemize}
Для б\'{о}льших длин берут модули больше машинного слова (тогда сами вычеты "--- длинные числа,
а умножение на $2^k$ остаётся сдвигом; на этом построен алгоритм Шёнхаге"--~Штрассена,
работающий в кольцах $\Zn{(2^N + 1)}$~\cite{schonhage-strassen}), либо несколько простых модулей
вида $c \cdot 2^k + 1$ с последующим восстановлением результата по КТО.

Итог: кольца Мерсенна и Ферма привлекательны тем, что в них
(а) приведение по модулю "--- одно сложение или вычитание половин числа;
(б) корень из единицы "--- $\pm 2$ или степень двойки, и умножение на него "--- сдвиг;
(в) свёртка вычисляется точно, без ошибок округления; а платой служат ограничения на длину
преобразования и неудобное представление (вычет $2^n$ в кольце Ферма).

\checkyourself{Почему ТЧП Мерсенна с $\omega = 2$ имеет простую длину $q$ и почему это плохо
для БПФ? Как получить длину $2q$?}

\checkyourself{Сравните методы Монтгомери, Барретта и специальные модули: какой из них
применим к модулю RSA, а какой "--- к модулю, выбираемому разработчиком протокола?}

\checkyourselfsection

\begin{furtherreading}
  \readitem[разд.~9.2.1, 9.2.3]{crandall-pomerance}{метод Монтгомери, модули специального
    вида; разд.~9.5 "--- свёртки и теоретико-числовые преобразования}
  \readitem[гл.~1]{chernov-korepanov}{теоретико-числовые преобразования, преобразования
    Мерсенна и Ферма, их применение для вычисления свёрток}
  \readitem[разд.~14.3]{hac}{приведение по модулю: классическое, Монтгомери, Барретта,
    модули специального вида (с примерами)}
  \readitem[гл.~2]{brent-zimmermann}{модулярная арифметика: редукции Барретта и Монтгомери,
    специальные модули, умножение в кольцах $\Zn{(2^N \pm 1)}$}
  \readitem{montgomery}{оригинальная статья Монтгомери (три страницы)}
  \readitem{rader-mersenne}{преобразование Мерсенна (оригинальная статья)}
  \readitem{agarwal-burrus}{преобразование Ферма и его применение к цифровой фильтрации}
  \readitem{leibowitz}{код с уменьшением на единицу для арифметики по модулю $2^n + 1$}
  \readitem{cp-algorithms}{статьи «Montgomery Multiplication» и «Fast Fourier transform»
    (раздел о вычислениях по модулю)}
\end{furtherreading}
\end{lecture}
